God names a king who does not know him. Jesus asks his opponents to name the image on a coin. Between those two acts of naming stands a question that reaches beyond taxation: how can a person receive the goods of this world without giving them the place that belongs to God? Cyrus has strength, a royal task and a place in Israel's deliverance. None of these makes him the source of the salvation he serves. Caesar has a recognizable currency and real claims within civil life. Neither the image nor the inscription makes his claim unlimited.

The Twenty-ninth Sunday's worship begins elsewhere than the court or the marketplace. It begins with a petitioner asking to be heard and sheltered. The Collect concerns the will that must serve; the prayers at the altar concern gifts, purification and a life moving toward eternal fulfilment. These prayers bring the argument home. The assembly is not merely deciding how powerful outsiders should behave. It must ask how its own powers, possessions, service and hope are to be returned to their giver.

Paul's thanksgiving adds a second question. The Thessalonians have received a Gospel that bears fruit in work, labour and endurance. How can their action be real if the gift comes first? Their faith does not make them spectators of God's power. Divine power is known in a community that believes, loves and perseveres. The short acclamation from Philippians holds together a received word and a visible light. The Eucharistic movement then carries this relation between gift and response into worship: those who offer still need cleansing, and those who receive still need to grow.

Two readings unfold these questions. The first follows Jerome's distinction between Cyrus and the true Christ, and Chrysostom's distinction between rightful dues and impious demands. It finds a freedom to serve the common life without worshipping its powers. The second follows Aquinas's account of God's action in willing and Chrysostom's account of faith made active under pressure. It finds a freedom that grows through receiving and answering grace. Each reaches the whole formulary, but their answers are different: the first orders our allegiances; the second examines the transformation of the person who owes them.

The Communion alternatives give that distinction a particularly concrete setting. One uses a Psalm of mercy, deliverance and nourishment. The other selects the Son of Man's self-gift as ransom. The former teaches dependence amid need; the latter names the saving deed on which Christian service rests. They are alternatives, not consecutive texts that every congregation hears. Either can deepen both readings, while changing the image through which their conclusion is received.
